
\section{Preliminaries}
\label{sect:prelim} 
Relational doctrines \cite{DagninoP23,DagninoP25tac} were introduced as a functorial description of (the core fragment of) the calculus of relations \cite{Tarski41}, 
that is, the fragment whose operations are only the identity, the composition and the converse. 
They can be synthetically defined as the faithful framed bicategories \cite{Shulman08} with a dagger or, more explicitly, as follows.

\begin{definition}[Relational Doctrine]\label[def]{def:rel-doc}
A \emph{relational doctrine} consists of a base category \CC and a functor \fun{\RDoc}{\bop\CC}{\Pos} such that for every triple of objects $X,Y,Z$ in \CC there is a monotone function \[\fun{\blank\rcomp\blank}{\RDoc(X,Y)\times\RDoc(Y,Z)}{\RDoc(X,Z)}\] and an element $\rid_X \in \RDoc(X,X)$ such that 

\begin{align*} 
\relr\rcomp(\rels\rcomp\relt) &= (\relr\rcomp\rels)\rcomp\relt 
& 
\rid_X\rcomp\relr &= \relr 
&
\relr\rcomp\rid_Y &= \relr 
\\
(\relr\rcomp\rels)\rconv &= \rels\rconv \rcomp \relr\rconv 
&
\rid_X\rconv &= \rid_X
&
\relr\rdconv &= \relr 
\end{align*} 
and such that for all $\relr\in\RDoc(X,Y)$, $\rels\in\RDoc(Y,Z)$ and 
\fun{f}{A}{X}, \fun{g}{B}{Y} and \fun{h}{C}{Z}  in \CC it holds that
\[\RDoc\reidx{f,g}(\relr)\rcomp\RDoc\reidx{g,h}(\rels) \order \RDoc\reidx{f,h}(\relr\rcomp\rels)
\qquad
\rid_X \order \RDoc\reidx{f,f}(\rid_Y)
\qquad
(\RDoc\reidx{f,g}(\relr))\rconv \order \RDoc\reidx{g,f}(\relr\rconv)\] 
\end{definition}


The element $\rid_X$ is the \emph{identity} or \emph{diagonal} relation on $X$, 
$\relr\rcomp\rels$  is the \emph{relational composition} of $\relr$ followed by $\rels$, and 
$\relr\rconv$ is the \emph{converse} of the relation $\relr$.  
% 
Note that all relational operations are lax natural transformations, but 
the operation of taking the converse, being  an involution, is actually strictly natural. 
% 
The requirement of lax naturality becomes clear if one looks at the examples. 
For instance, in the paradigmatic example of set-theoretic relations, given by the relational doctrine \fun{\Rel}{\bop\Set}{\Pos} 
where $\Rel(X,Y)=\PP(X\times Y)$ and  $\Rel(f,g)=(f\times g)^{-1}$; it is easy to check that 
the inclusion
$\Rel\reidx{f,g}(\relr)\rcomp\Rel\reidx{g,h}(\rels)
\subseteq\Rel\reidx{f,h}(\relr\rcomp\rels)$
is an equality if and only if $g$ is surjective and, similarly, the inclusion $\rid_X 
\subseteq \RDoc\reidx{f,f}(\rid_Y)$ is an equality if and only if $f$ is injective.

Some examples of relational doctrines can be found in \cite{DagninoP23,DagninoP25tac,DagninoP25apal}, here we just consider those that will play a relevant role in the paper.

\begin{example} \label[ex]{ex:rel-doc} 
\begin{enumerate}
% 
\item\label{ex:rel-doc:span}
Let \CC be a category with weak pullbacks. 
Denote by $\Span\CC(X,Y)$ the poset reflection of the preorder whose objects are spans in \CC between $X$ and $Y$ and where 
$\spn{X}{p_1}{A}{p_2}{Y} \order \spn{X}{q_1}{B}{q_2}{Y}$ if and only if there is an arrow \fun{f}{A}{B} such that 
$p_1 = q_1 \circ f $ and $p_2 = q_2 \circ f$. 
Given a span $\relr = \spn{X}{p_1}{A}{p_2}{Y}$ and arrows \fun{f}{X'}{X} and \fun{g}{Y'}{Y} in \CC, define 
$\Span\CC\reidx{f,g}(\relr) \in \Span\CC(X',Y')$ by one of the following equivalent diagrams: 
\[
\vcenter{\xymatrix@C=3ex@R=3ex{
&& W \ar[ld] \ar[rrdd] \ar@{}[rddd]|{wpb} && \\
& W'\ar[ld] \ar[rd] \ar@{}[dd]|{wpb} &&& \\  
X' \ar[rd]_-{f} && A \ar[ld]^-{p_1} \ar[rd]_-{p_2} && Y' \ar[ld]^-{g} \\ 
& X && Y 
}} 
\quad 
\vcenter{\xymatrix@C=3ex@R=3ex{
&& W \ar[rd] \ar[lldd] \ar@{}[lddd]|{wpb} && \\
&&& W'\ar[rd] \ar[ld] \ar@{}[dd]|{wpb} & \\  
X' \ar[rd]_-{f} && A \ar[ld]^-{p_1} \ar[rd]_-{p_2} && Y' \ar[ld]^-{g} \\ 
& X && Y 
}} 
\]
The functor \fun{\Span\CC}{\bop\CC}{\Pos} is a relational doctrine where, for 
$\relr = \spn{X}{p_1}{A}{p_2}{Y}$ and 
$\rels = \spn{Y}{q_1}{B}{q_2}{Z}$ it is 
\[
\rid_X = \vcenter{\xymatrix@R=3ex@C=3ex{
& X \ar[ld]_-{\id_X} \ar[rd]^-{\id_X} & \\ 
X && X
}} \qquad 
\relr \rcomp \rels = \vcenter{\xymatrix@R=3ex@C=3ex{
&& W \ar[ld] \ar[rd] \ar@{}[dd]|{wpb} && \\ 
& A \ar[ld]^-{p_1}\ar[rd]_-{p_2} && B \ar[ld]^-{q_1} \ar[rd]_-{q_2} & \\ 
X && Y && Z 
}}
\]
while $\relr\rconv=  \spn{Y}{p_2}{A}{p_1}{X}$% 
\item\label{ex:rel-doc:jmspans}
Suppose \CC is regular  and 
denote by $\jmSpan\CC(X,Y)$ the subposet of $\Span\CC(X,Y)$ on jointly monic spans. This extends to a functor  \fun{\jmSpan\CC}{\bop\CC}{\Pos} which is a relational doctrine, where the identity relations and the converses are defined as in \cref{ex:rel-doc:span}, while to compute the relational composition of two jointly monic spans $\relr$ and $\rels$ it suffices to observe that any spans factors through a jointly monic one.
% 
\item\label{ex:rel-doc:vrel}
Let $\RPos = \ple{[0,\infty],\ge,+,0}$ be the quantale of extended non-negative real numbers  (here addition and subtraction in $[0,\infty]$ are extensions of those in $[0,\infty)$ by the rules $\infty+a=a+\infty=\infty-a=\infty$ and $a-\infty=0$ see \cite{Lawvere73}).
A \emph{metric relation} between sets $X$ and $Y$ is just a function \fun{\alpha}{X\times Y}{[0,\infty]}. 
The functor \fun{\VRel\RPos}{\bop\Set}{\Pos} maps two sets $X$ and $Y$ to 
$\VRel\RPos(X,Y) = [0,\infty]^{X\times Y}$, the set of metric relations from $X$ to $Y$  ordered pointwise, while $\VRel\RPos\reidx{f,g}$ is the precomposition with $f\times g$. 
The identity relation, the relational composition and the converse of a relation are defined as follows: 
\begin{align*}
\rid_X(x,x') &= \begin{cases}
0 & x = x' \\
\infty & x \ne x' 
\end{cases}
\\ 
(\relr\rcomp\rels)(x,z) &= \inf_{y\in\Car Y} (\relr(x,y) + \rels(y,z)) 
\\ 
\relr\rconv(y,x) &= \relr(x,y)
\end{align*}
%
where $\relr\in\VRel\RPos(X,Y)$ and $\rels\in\VRel\RPos(Y,Z)$. 

\item\label{ex:rel-doc:met} 
 Following Lawvere~\cite{Lawvere73}, a (symmetric) metric space $X$ consists of a set $\Car{X}$ and a function 
$\fun{\dist_X}{\Car{X}\times\Car{X}}{[0,\infty]}$ such that $\dist_X(x,x)=0$ and $\dist_{X}(x,x')=\dist_{X}(x',x)$ and $\dist_{X}(x,x')+\dist_{X}(x',x'')\ge\dist_{X}(x,x'')$. 
A metric space $X$ is \emph{separated} if additionally $\dist_X(x,x') = 0$ implies $x = x'$. 
Thus, separated metric spaces are essentially usual metric spaces where we also admit points with an infinite distance. 
Let $\Met$ be the category of separated metric spaces and non-expansive maps, i.e., 
arrows $\fun{f}{X}{Y}$ are functions $\fun{f}{\Car{X}}{\Car{Y}}$ such that 
$\dist_X(x,x') \ge \dist_Y(f(x),f(x'))$. 
A \emph{bimodule} $\phi$ between $X$ and $Y$ 
is a function 
from $\Car{X}\times\Car{Y}$ such that 
$\dist_{X}(x,x')+\phi(x,y)+\dist_{Y}(y,y')\ge\phi(x',y')$. 
We denote by $\BM(X,Y)$ the poset of bimodules between $X$ and $Y$ with the pointwise order. 
This assignment extends to a functor  $\fun{\BM}{\bop\Met}{\Pos}$ that acts on pairs of arrows by precomposition. 
The functor $\BM$ is a relational doctrine, where relational composition and the converse of a relation are defined as in \cref{ex:rel-doc:vrel} and where the identity $\rid_X$ is $\dist_X$.  
\end{enumerate}
\end{example}

\begin{remark} 
The relational doctrines presented in \cref{ex:rel-doc:vrel,ex:rel-doc:met} of \cref{ex:rel-doc} belong to a larger class of relational doctrines, 
obtained by replacing $\RPos$ with an arbitrary commutative quantale. 
All the results we will present for these specific case easily extend to this larger class. 
\end{remark}

Relational doctrines are the objects of the 2-category \RDtn. 

A 1-arrow \oneAr{F}{\RDoc}{\SDoc} from \fun{\RDoc}{\bop\CC}{\Pos} to \fun{\SDoc}{\bop\D}{\Pos}, 
is a pair \ple{\fn{F},\lift{F}} consisting of 
a functor \fun{\fn{F}}{\CC}{\D} and a natural transformation 
\nt{\lift{F}}{\RDoc}{\SDoc\circ \bop{\fn{F}}},

\[
\xymatrix@C=7.5em@R=1em{
{\bop\CC}\ar[rd]^(.4){\RDoc}_(.4){}="P"
\ar[dd]_{\bop{\fn{F}}}^{}="F"
&\\
 & {\ct{Pos}}\\
{\bop\D}\ar[ru]_(.4){\SDoc}^(.4){}="R"&
\ar"P";"R"_{\lift{F}\kern.5ex\cdot\kern-.5ex}="b"
}
\]
preserving relational identities, composition and converse, that is 
\begin{align*}
\rid_{\fn{F}X} &= \lift{F}_{X,X}(\rid_X)\\ 
\lift{F}_{X,Y}(\relr)\rcomp \lift{F}_{Y,Z}(\rels) &= \lift{F}_{X,Z}(\relr\rcomp\rels)\\ 
\lift{F}_{X,Y}(\relr))\rconv &= \lift{F}_{Y,X}(\relr\rconv)
\end{align*}
%
for $\relr\in\RDoc(X,Y)$ and $\rels\in\RDoc(Y,Z)$. 

A 2-arrow \twoAr{\theta}{F}{G} is a natural transformation \nt{\theta}{\fn{F}}{\fn{G}} such that 
$\lift{F}_{X,Y} \order \SDoc\reidx{\theta_X,\theta_Y}\circ \lift{G}_{X,Y}$, for all objects $X,Y$ in the base of $\RDoc$
\[
\xymatrix@C=18em@R=1.5em{
{\bop\CC}\ar[rd]^(.4){\RDoc}_(.4){}="P"
\ar@<-1ex>@/_/[dd]_{\bop{\fn{F}}}^{}="F"\ar@<1ex>@/^/[dd]^{\bop{\fn{F'}}}_{}="G"&\\
 & {\ct{Pos}}\\
{\bop\D}\ar[ru]_(.4){\SDoc}^(.4){}="R"&
\ar@/_/"P";"R"_{\lift{F}\kern.5ex\cdot\kern-.5ex}="b"
\ar@<1ex>@/^/"P";"R"^{\kern-.5ex\cdot\kern.5ex \lift{F'}}="c"
\ar"G";"F"_{.}^{\theta\op}\ar@{}"b";"c"|{\le}}
\]
 % 

We briefly recall some basic facts on relational doctrines and refer the reader to  \cite{DagninoP23,DagninoP25apal} for details. 
Fix a relational doctrine \fun{\RDoc}{\bop{\CC}}{\Pos}. 

For every arrow \fun{f}{X}{Y} in \CC the relation $\gr{f}=\RDoc\reidx{f,\id_Y}(\rid_Y)\in\RDoc(X,Y)$ is called \emph{graph} of $f$. 
Graphs commutes with composition and identities in the sense that  
$\gr{g\circ f} = \gr{f}\rcomp\gr{g}$ and $\gr{\id_X} = \rid_X$. 

An immediate use of graphs is in writing reindexing \fun{\RDoc\reidx{f,g}}{\RDoc(X,Y)}{\RDoc(A,B)} in relational terms and, moreover, in showing that each $\RDoc\reidx{f,g}$ has a left adjoint \fun{\Ex^\RDoc\reidx{f,g}}{\RDoc(A,B)}{\RDoc(X,Y)}. Indeed 
\[
\RDoc\reidx{f,g}(\relr) = \gr{f} \rcomp \relr \rcomp \gr{g}\rconv 
\qquad 
\Ex^\RDoc\reidx{f,g}(\rels) = \gr{f}\rconv \rcomp \rels \rcomp \gr{g} 
\]
%
The condition $\gr{f}=\gr{g}$ (or equivalently  $\rid_X\le \gr{f}\rcomp\gr{g}\rconv$ or also $\rid_X\le \gr{g}\rcomp\gr{f}\rconv$) defines a notion of equality between the arrows of $\CC$ which is, in general, weaker then the actual equality (see \cite{DagninoP25tac,DagninoP25apal} for a more detailed analysis). 
% 
We say that  relational doctrine $\fun{\RDoc}{\bop\CC}{\Pos}$ is \emph{extensional} when for all \fun{f,g}{X}{Y} in \CC, if $\gr{f} = \gr{g}$, then $f = g$. 

\begin{example}\label[ex]{ex:rel-doc-ex}
All relational doctrines presented in \cref{ex:rel-doc} are extensional. In particular doctrines of the form $\Span\CC$ and $\jmSpan\CC$ in \refItem{ex:rel-doc}{span} and in \refItem{ex:rel-doc}{jmspans} are extensional as in both examples $\rid_X$ is $\spn{X}{\id_X}{X}{\id_X}{X}$; the relational doctrines $\VRel\RPos$ presented in \refItem{ex:rel-doc}{vrel} is extensional as for every set $X$ the metric relation $\rid_X$ is two-valued; finally,  the relational doctrines $\BM$ presented in \refItem{ex:rel-doc}{met} is extensional because metric spaces in $\Met$ are separated: 
if $\dist_Y(f(x),y) = \gr{f}(x,y) = \gr{g}(x,y) = \dist_Y(g(x),y)$, 
then $0=\dist_Y(f(x),f(x))=\dist_Y(g(x),f(x))$ so $f(x)=g(x)$.
\end{example}


A relation  $\relr\in\RDoc(X,Y)$ is
\begin{center}
\begin{tabular}{ll}
\emph{functional} if $\relr\rconv \rcomp \relr \order \rid_Y$&\emph{injective} if $\relr\rcomp\relr\rconv \order \rid_X$\\[1ex]
\emph{total} if $\rid_X \order \relr\rcomp\relr\rconv$&\emph{surjective} if $\rid_Y \order \relr\rconv \rcomp \relr$
\end{tabular}
\end{center}

For every $f$, the graph $\gr{f}$ is functional and total. The next proposition shows that 
 functional and total relations are discretely ordered.

\begin{proposition}\label[prop]{prop:fun-ord}
Let $\RDoc$ be a relational doctrine over \CC. 
For functional and total relations $\relr,\rels\in\RDoc(X,Y)$
if $\relr\order\rels$, then $\relr = \rels$. 
\end{proposition}
\begin{proof} 
If $\relr\order\rels$, then 
$\rels = \rid_X\rcomp\rels 
       \order \relr \rcomp \relr\rconv \rcomp \rels 
       \order \relr \rcomp \rels\rconv \rcomp \rels 
       \order \relr \rcomp \rid_Y 
       = \relr$.
\end{proof}


We will say that an arrow $f$ is 
\emph{$\RDoc$-injective} (resp. \emph{$\RDoc$-surjective}) or simply injective when $\RDoc$ is clear, when its graph $\gr{f}$ is injective (resp. surjective). A arrow is \emph{$\RDoc$-bijective} when it is  $\RDoc$-injective and $\RDoc$-surjective.

In general, monomorphisms (resp. epimorphisms) of $\CC$ and $\RDoc$-injections (resp. $\RDoc$-surjections) determines different classes of arrows, unless one assumes some extra structure such as in the following two propositions.

\begin{proposition}\label{prop:extomonic}In every extensional relational doctrine \fun{\RDoc}{\bop\CC}{\Pos} each $\RDoc$-injective arrow is monic, each $\RDoc$-surjective arrow is epic.
\end{proposition}
\begin{proof} Take a $\RDoc$-injective arrow $\fun{m}{A}{B}$ and let $\fun{f,g}{X}{B}$ be such that $mf=mg$. Then $\gr{mf}=\gr{mg}$ so $\rid_X\le \gr{f}\rcomp\gr{m}\rcomp\gr{m}\rconv\rcomp\gr{g}\rconv$. By $\RDoc$-injectivity of $m$ one has $\rid_X\le \gr{f}\rcomp\gr{g}\rconv$, then $f=g$ by extensionality. A similar argument proves that $\RDoc$-surjective arrow are epimorphisms.
\end{proof}

\begin{proposition}\label{prop:splitepi}
In every  relational doctrine \fun{\RDoc}{\bop\CC}{\Pos} retractions in $\CC$ are $\RDoc$-surjective and sections are $\RDoc$-injective.
\end{proposition}
\begin{proof} 
Let $\fun{e}{X}{Y}$ and $\fun{r}{Y}{X}$ be such that  $\id_Y=er$, then $\rid_Y=\gr{\id_Y}=\gr{er}=\gr{r}\rcomp\gr{e}$. So $\rid_Y= \rid_Y\rconv\rcomp\rid_Y=\gr{e}\rconv\gr{r}\rconv\rcomp\gr{r}\rcomp\gr{e}$ and since $\gr{r}$ is functional $\rid_Y\le \gr{e}\rconv\rcomp\gr{e}$. Similarly one proves the $\RDoc$-injectivity of $r$. 
\end{proof}




\subsection{Quotients and the extensional quotient completion}
Let $X$ be an object of \CC. 
An \emph{$\RDoc$-equivalence relation} on $X$ is a relation 
$\eqrelr$ in $\RDoc(X,X)$ satisfying
\begin{center}
\begin{tabular}{lll}
reflexivity: $\rid_X\order\eqrelr$& symmetry: $\eqrelr\rconv\order\eqrelr$&transitivity: $\eqrelr\rcomp\eqrelr\order\eqrelr$
\end{tabular}
\end{center}
Every arrow \fun{f}{X}{Y} in \CC induces a $\RDoc$-equivalence relation on $X$, dubbed \emph{kernel} of $f$, given by $\gr{f}\rcomp\gr{f}\rconv$.
One immediately sees that for the relational doctrine \fun{\Rel}{\bop\Set}{\Pos} of set-theoretic relations the kernel of a function $f$ is precisely the set $\{(x,x')\mid f(x)=f(x')\}$. 

A \emph{quotient arrow} of $\eqrelr$ is an arrow \fun{q}{X}{W} in \CC such that 
$\eqrelr \order \gr{q}\rcomp\gr{q}\rconv$ and, 
such that for every arrow \fun{f}{X}{Z} with $\eqrelr\order\gr{f}\rcomp\gr{f}\rconv$, there is a unique arrow \fun{h}{W}{Z} with $f = h\circ q$.  
We say that a quotient arrow \fun{q}{X}{W}  for $\eqrelr$ is 
\emph{effective} if $\eqrelr = \gr{q}\rcomp\gr{q}\rconv$, i.e., its kernel coincides with the equivalence relation, 
and it is \emph{descent} if $\gr{q}\rconv\rcomp\gr{q} = \rid_W$, i.e., $q$ is sujective. 
%
Finally, we say that \emph{$\RDoc$ has quotients} if every $\RDoc$-equivalence relation admits an effective descent quotient arrow. 


\begin{example} \label[ex]{ex:rel-doc-q} 
\begin{enumerate}
% 
\item\label{ex:rel-doc-q:span}
For the relational doctrine $\Span\CC$ of spans over a category \CC with weak pullbacks (see \refItem{ex:rel-doc}{span})
a $\Span\CC$-equivalence relation is a pseudo-equivalence relation of \CC in the sense of \cite{CarboniA:regec}.
% 
\item\label{ex:rel-doc-q:jmspans}
For the relational doctrine  $\jmSpan\CC$ of jointly monic spans over a regular category $\CC$ (see \refItem{ex:rel-doc}{jmspans}) a $\jmSpan\CC$-equivalence relation  is an equivalence relation in \CC and quotient arrows are regular epimorphisms. 
Moreover  \CC is exact if and only if $\jmSpan\CC$ has quotients. 
% 
\item\label{ex:rel-doc-q:vrel}
For the relational doctrine $\VRel\RPos$ of metric relations  (see  \refItem{ex:rel-doc}{vrel}) a $\VRel\RPos$-equivalence relation $d$ over a set $X$ is a function $d:X\times X\to [0,\infty]$ suchthat $d(x,x')=0$ (reflexivity), $d(x,x')=d(x',x)$ (symmetry) and $d(x,x')+d(x',x'')\geq d(x,x'')$ (transitivity); in other words, the $\VRel\RPos$-equivalence relations over a set $X$ are the distances on $X$.
% 
\item\label{ex:rel-doc-q:met} 
Consider the relational doctrine  $\fun{\BM}{\bop\Met}{\Pos}$ of bimodules over separated metric spaces as in \refItem{ex:rel-doc}{met}. 
A $\BM$-equivalence relation over a metric space $X$ is a distance $\rho$ over $\Car{X}$ such that $\dist_{X}(x,x')\geq \rho(x,x')$. 
Note that these metrics are not required to be separated.  Therefore, in order to build a quotient, we have to force separation. 
To this end, given a $\BM$-equivalence relation $\rho$ over $X$, we consider the equivalence relation $\sim_\rho$ on the set $\Car{X}$ defined by $x\sim_\rho x'$ if and only if $\rho(x,x') = 0$ and the metric space 
$X_\rho$ where $\Car{X_\rho} = \Car{X}/\sim_\rho$ and $\dist_{X_\rho}(\eqc{x},\eqc{x'}) = \rho(x,x')$. 
It is easy to check that $X_\rho$ is a separated metric space. 
Moreover, the quotient projection function $\fun{\pi_\rho}{\Car{X}}{\Car{X}/\sim_\rho}$ is non-expansive and gives a quotient arrow $\fun{\pi_\rho}{X}{X_\rho}$ of $\rho$ in $\BM$. 
More generally, one can prove that an arrow in $\Met$ is a quotient arrow for a $\BM$-equivalence relation if and only if its underlying function is surjective. 
The underling function of $\pi_\rho$ is trivially surjective. 
Conversely, given an arow $\fun{f}{X}{Y}$ whose underlying function is surjective, 
denote by $k_f$ the kernel of $f$, that is, the $\BM$-equivalence relation on $X$ defined by 
$k_f(x,x') = \dist_Y(f(x),f(x'))$ 
and take the quotient projection $\fun{\pi_{k_f}}{X}{X_{k_f}}$. 
The function $\fun{j}{\Car{X_{k_f}}}{\Car{Y}}$, determined by $j(\eqc{x}) = f(x)$ is well-defined, since $Y$ is separated, and non-expansive. 
It is an isomorphism such that $j\circ \pi_{k_f}=f$, proving that $f$ is a quotient arrow. 
\end{enumerate}
\end{example}

Any 1-arrow \oneAr{F}{\RDoc}{\SDoc}  in \RDtn 
preserves equivalence relations, that is, 
if $\eqrelr$ is an $\RDoc$-equivalnece relation on $X$, then $\lift{F}_{X,X}(\eqrelr)$ is an $\SDoc$-equivalence relation on $\fn{F}X$. 
% 
We say that a 1-arrow \oneAr{F}{\RDoc}{\SDoc} preserves quotients if, 
whenever $q$ is a quotient arrow for an $\RDoc$-equivalence relation $\eqrelr$ on $X$, the arrow $\fn{F}q$ is a quotient arrow for the $\SDoc$-equivalence relation $\lift{F}_{X,X}(\eqrelr)$ on $\fn{F}X$. 
% 
We denote by $\QRDtn$ the 2-full 2-subcategory of $\RDtn$ whose objects are relational doctrines with quotients and whose 1-arrows are those of $\RDtn$ that preserves quotients. 

We now describe the construction presented in \cite{DagninoP23, DagninoP25apal}, that takes a relational doctrine 
\fun{\RDoc}{\bop\CC}{\Pos} and produces another relational doctrine 
\fun{\EQR\RDoc}{\bop{\EQC\RDoc}}{\Pos} which is extensional and has quotients. 
% 
The category $\EQC\RDoc$ is defined as follows: 
\begin{itemize}
\item an object is a pair \ple{X,\eqrelr} of an object $X$ in \CC and an $\RDoc$-equivalence relation on it; 
\item an arrow \fun{\eqc{f}}{\ple{X,\eqrelr}}{\ple{Y,\eqrels}} is an equivalence class of arrows \fun{f}{X}{Y} in \CC such that $\eqrelr \order \gr{f}\rcomp\eqrels\rcomp\gr{f}\rconv$, modulo the equivalence reltion stating that two such arrow $f$ and $g$ are equivalent if and only if $\eqrelr\order\gr{f}\rcomp\eqrels\rcomp\gr{g}\rconv$. 
\end{itemize}
The functor $\EQR\RDoc$ is then defined as follows:
\[
\EQR\RDoc(\ple{X,\eqrelr},\ple{Y,\eqrels}) = \{ \relr\in\RDoc(X,Y) \mid \eqrelr\rcomp\relr\rcomp\eqrels\order\relr \}
\qquad 
\EQR\RDoc\reidx{\eqc{f},\eqc{g}} = \RDoc\reidx{f,g} 
\]
The elements of the fibre $\EQR\RDoc(\ple{X,\eqrelr},\ple{Y,\eqrels})$ are called \emph{descent data} for $\eqrelr$ and $\eqrels$. 
% 
Finally, the relational operations of $\EQR\RDoc$ are the same as in $\RDoc$ except the identity which is given by $\rid_{\ple{X,\eqrelr}} = \eqrelr$. 

The relational doctrine $\EQR\RDoc$ has quotients: 
an $\EQR\RDoc$-equivalence relation $\eqrels$ on $\ple{X,\eqrelr}$ is  an $\RDoc$-equivalence relation on $X$ such that $\eqrelr\order\eqrels$ and the quotient arrow of $\eqrels$ is \fun{\eqc{\id_X}}{\ple{X,\eqrelr}}{\ple{X,\eqrels}}. 



\begin{example} \label[ex]{ex:rel-doc-eqc} 
\begin{enumerate}
\item\label{ex:rel-doc-eqc:span}
If \CC has weak finite limits, 
then $\EQR{\Span\CC}$ is isomorphic to the doctrine of jointly monic spans over  $\CC_{ex/wlex}$, the exact completion of a weakly lex category as in \cite{CarboniA:regec} (crf. also \cref{rem:MillyPino}).  We did not find a reference for this claim whose proof is a bit lengthy and, as such, it is postponed to \cref{ss:spans}. It can be safetly skipped by those readers who are not interested in it.

\item\label{ex:rel-doc-eqc:vrel}
As proved in \cite{DagninoP23,DagninoP25apal}, the relational doctrine $\EQR{\VRel\RPos}$ is equivalent in $\EQRDtn$ to the relational doctrine $\fun{\BM}{\bop{\Met}}{\Pos}$ of bimodules over  separated metric spaces.
\end{enumerate}
\end{example}


\begin{remark}\label{rem:MillyPino} 
In \cite{DagninoP25apal} we showed that relational doctrines that are cartesian (i.e the 1-arrows $\fun{!_\RDoc}{\RDoc}{1}$ and $\fun{\Delta_\RDoc}{\RDoc}{\RDoc\times \RDoc}$
have a right adjoint in \RDtn) and modular (i.e. relations satisfy $\alpha\rcomp\beta \wedge \gamma \le \alpha\rcomp(\beta\wedge \alpha\rconv\rcomp\gamma)$) are equivalent to elementary and existential doctrines as in \cite{MaiettiME:quofcm}. 
% 
If $\RDoc$ is cartesian and modular, so is $\EQR{\RDoc}$ and, in this case, $\EQR{\RDoc}$ coincides with the completion introduced by Maietti and Rosolini in \cite{MaiettiME:eleqc,MaiettiME:quofcm}, dubbed elementary quotient completion. So in particular the examples listed in \cite{MaiettiME:eleqc,MaiettiME:quofcm} such as the $ex/wlex$ completion of a category $\CC$ with strong products, and other non-exact categories of equivalence relations, such as the category of equilogical spaces, the category of assemblies over a partial combinatory algebra, the categories of $H$-valued partial equivalence relations for a fixed locale $H$ and categories of setoids over type theories are examples of relational quotient completion for an appropriately chosen relational doctrine.
\end{remark}


We conclude the section recalling that the construction of $\EQR\RDoc$ is universal in the sense of the following proposition, whose proof can be found in \cite{DagninoP25apal}. Let $\EQRDtn$ be the full 2-subcategory of $\QRDtn$ on those relational doctrines with quotients which are also extensional, denote by $\fun{\EQRFun}{\EQRDtn}{\RDtn}$ the forgetful functor and observe that the relational quotient completion of a relational doctrine, i.e. the construction $\RDoc\mapsto\EQR\RDoc$, extends to a 2-functor $\fun{\REQFun}{\RDtn}{\EQRDtn}$.

\begin{proposition}\label{prop:eqc-mnd}The relational quotient completion determines a biadjunction $\REQFun\dashv\EQRFun$. 
\end{proposition}


\subsection{Quotient completion of doctrines of spans}\label{ss:spans}

Here we prove the claim in \refItem{ex:rel-doc-eqc}{span} stating that, for a category $\CC$ with weak finite limits, $\EQR{\Span\CC}$ is isomorphic in $\EQRDtn$ to $\fun{\jmSpan{\CC_{ex/wlex}}}{\bop{(\CC_{ex/wlex})}}{\Pos}$. 
As already mentioned, we are including this full proof for sake of completeness, as we have not found it in the literature. 
However, it is not essential for understanding the rest of the paper.



First of all, let us notice that 
the base of $\EQR{\Span\CC}$ is precisely $\CC_{ex/wlex}$. 
To show that also the doctrines are isomorphic we  need some instrumental lemmas.

\begin{lemma}\label{l:isofrecce}
Suppose $\RDoc$ and $\SDoc$ are relational doctrines over the same base $\CC$ and \oneAr{F}{\RDoc}{\SDoc} is a 1-arrow where $\fun{\fn{F}}{\CC}{\CC}$ is the identity on $\CC$. If for every $X,Y$ in $\CC$ the map $\fun{\lift{F}_{X,Y}}{\RDoc(X,Y)}{\SDoc(X,Y)}$ is an isomorphism in $\Pos$, then $\RDoc$ and $\SDoc$ are isomorphic.
\end{lemma}
\begin{proof}For every $X,Y$, call $G_{X,Y}$ the inverse of $\lift{F}_{X,Y}$ in $\Pos$.  The following equations 
\begin{align*} 
G_{X,X}(\rid_X)&=G_{X,X}(\lift{F}_{X,X}(\rid_X))=\rid_X\\
G_{X,Z}(\relr\rcomp\rels)&=G_{X,Z}(\lift{F}_{X,Y}G_{X,Y}(\relr)\rcomp \lift{F}_{Y,Z}G_{Y,Z}(\rels))\\
&=G_{X,Z}\lift{F}_{X,Z}(G_{X,Y}(\relr)\rcomp G_{Y,Z}(\rels))=G_{X,Y}(\relr)\rcomp G_{Y,Z}(\rels)\\
G_{X,Y}(\relr\rconv)&= G_{X,Y}((\lift{F}_{X,Y}G_{X,Y}(\relr))\rconv)=G_{X,Y}\lift{F}_{X,Y}((G_{X,Y}(\relr))\rconv)=G_{X,Y}(\relr)\rconv
\end{align*} 
 show that the inverses commute with the relational operations, as such they determine a 1-arrow of relational doctrines $\fun{}{\SDoc}{\RDoc}$.
\end{proof}

Let $\ct{E}$ be an exact category and $\ct{P}$ a regular projective cover of it. Denote by $\Span{\ct{E}}_{\ct{P}}(X,Y)$ the subposet of $\Span{\ct{E}}(X,Y)$ on those spans whose vertex is in $\ct{P}$. If $\relr=\spn{X}{f}{A}{g}{Y}$ is a span  in $\Span{\ct{E}}(X,Y)$ and $\fun{q}{P}{A}$ and $\fun{q'}{P'}{A}$ are regular epimorphisms, then,  $q$ and $q'$ factors through each other, so the spans $\spn{X}{fq}{P}{gq}{Y}$ and $\spn{X}{fq'}{P'}{gq'}{Y}$ represent the same element in $\Span{\ct{E}}_{\ct{P}}(X,Y)$. 
We denote it by $\lift{\coveredarr}(\relr)$.



Note also that a regular projective cover of a pullback in $\ct{E}$ determines a weak pullback, i.e.  given $\fun{e}{X}{Y}$ and $\fun{f}{A}{Y}$ the squares $f\circ (f^*e\circ q)=e\circ (e^*f\circ q)$, where $q$ is a regular epimorphism with a regular projective domain, is a weak pullback. 
As a consequence of this and since regular epimorphisms compose, for every $f$ and $e$ as above there is a weak pullback of $e$ along $f$ 
which is a regular epimorphism whose domain is in $\ct{P}$.
This suffices to show that (equivalence classses of) spans with a regular projective vertex form a relational doctrine $\fun{\Span{\ct{E}}_{\ct{P}}}{\bop{\ct{E}}}{\Pos}$ where the composition of $\relr$ and $\rels$ is $\lift{\coveredarr}(\relr\rcomp\rels)$ and the identity relation over $A$ is $\lift{\coveredarr}(\rid_A)$ and $(\Span{\ct{E}}_{\ct{P}})_{f,g}= \lift{\coveredarr}(\Span{\ct{E}}_{f,g})$. 
Moreover the way we defined the relational operations of $\Span{\ct{E}}_{\ct{P}}$ ensures that \oneAr{\coveredarr}{\jmSpan{\ct{E}}}{\Span{\ct{E}}_{\ct{P}}} is indeed a 1-arrow as depicted below 
\[
\xymatrix@C=7.5em@R=1em{
{\bop{\ct{E}}}\ar[rd]^(.4){\jmSpan{\ct{E}}}_(.4){}="P"
\ar[dd]_{\bop{\fn{\coveredarr}}}^{}="F"
&\\
 & {\ct{Pos}}\\
{\bop{\ct{E}}}\ar[ru]_(.4){\Span{\ct{E}}_{\ct{P}}}^(.4){}="R"&
\ar"P";"R"_{\lift{\coveredarr}\kern.5ex\cdot\kern-.5ex}="b"
}
\]
 where \fun{\fn{\coveredarr}}{\ct{E}}{\ct{E}} is the identity functor on $\ct{E}$.

The following lemma is a stronger version of Lemma 35 in \cite{CarboniA:regec}. 

\begin{lemma} \label{l:analogo0} 
The 1-arrow \oneAr{\coveredarr}{\jmSpan{\ct{E}}}{\Span{\ct{E}}_{\ct{P}}} 
is an isomorphism.
\end{lemma}
\begin{proof} 
After \cref{l:isofrecce}, we only need to show that $\fun{\lift{\coveredarr}_{X,Y}}{\jmSpan{\ct{E}}(X,Y)}{\Span{\ct{E}}_{\ct{P}}(X,Y)}$ has an inverse in $\Pos$. This is given noticing that in an exact category every span factors through a jointly monic one, where the mediating arrow is a regular epimorphism.
\end{proof}

Let \CC be a category with weak finite limits and consider the relational doctrine $\fun{\Span\CC}{\bop\CC}{\Pos}$. 
Denote by  $\ct{C}_\Delta$ the regular projective cover of $\CC_{ex/wlex}$ determined by the objects of the form $(A,\rid_A)$. The following is an obvious corollary of \cref{l:analogo0} 

\begin{corollary}\label{c:analogo}
The doctrines $\Span{\ct{C_{ex/wlex}}}_{\ct{C}_\Delta}$ and $\jmSpan{\ct{C_{ex/wlex}}}$ are isomorphic.
\end{corollary}

The final stage of the proof consists in showing that $\Span{\ct{C_{ex/wlex}}}_{\ct{C}_\Delta}$ and $\EQR{\Span\CC}$ are isomorphic.
There is a morphism of doctrines \oneAr{i}{\EQR{\Span\CC}}{\Span{\ct{C_{ex/wlex}}}_{\ct{C}_\Delta}} where $\fn{i}$ is the identity functor on $\ct{C_{ex/wlex}}$ and the component of $\lift{i}$ at $\ple{X,\eqrelr}$ and $\ple{Y,\eqrels}$ sends a span $\spn{X}{f}{A}{g}{Y}$ to $\spn{\ple{X,\eqrelr}}{[f]}{\ple{A,\rid_A}}{[g]}{\ple{Y,\eqrels}}$. 

To build the inverse of $i$ the following lemma will be useful.

\begin{lemma} \label{l:chiudo} 
Let $\RDoc$ be a relational doctrine over $\CC$. For every equivalence relations $\eqrelr, \eqrels$ over $X$ and $Y$, respectively, the assignment $\relr\mapsto \eqrelr\rcomp\relr\rcomp\eqrels$ determines a monotone function $\fun{}{\RDoc(X,Y)}{\EQR\RDoc(\ple{X,\eqrelr},\ple{Y,\eqrels})}$.
\end{lemma}
\begin{proof} 
Monotonicity follows from monotonicity of relational composition. Recall that $\EQR\RDoc(\ple{X,\eqrelr},\ple{Y,\eqrels})= \{ \relr\in\RDoc(X,Y) \mid \eqrelr\rcomp\relr\rcomp\eqrels\order\relr \}$, thus $\eqrelr\rcomp\relr\rcomp\eqrels$ belongs to it by transitivity of $\eqrelr$ and $\eqrels$.
\end{proof}

Observe that an element  in $\Span{\CC}(X,Y)$, represented by $\spn{X}{f}{A}{g}{Y}$, can be written as $\gr{f}\rconv\rcomp\gr{g}$ where the relational operations are those of $\Span{\CC}$ and thus $\gr{f}$ is represented by the span 
$\spn{A}{\id_A}{A}{f}{X}$ and similarly for $\gr{g}$. 
Moreover, any span $\spn{\ple{X,\eqrelr}}{[f]}{\ple{A,\rid_A}}{[g]}{\ple{Y,\eqrels}}$ 
in $\ct{C_{ex/wlex}}$ is composed by legs, whose representative $f$ and $g$ are arrows of $\CC$ such that $\rid_A\order\gr{f}\rcomp\eqrelr\rcomp\gr{f}\rconv$ and $\rid_A\order\gr{g}\rcomp\eqrels\rcomp\gr{g}\rconv$ in $\Span{\CC}(A,A)$.

Thus, a candidate to be the inverse \oneAr{i'}{\Span{\ct{C_{ex/wlex}}}_{\ct{C}_\Delta}}{\EQR{\Span\CC}} 
of $i$ can be find choosing, for each $\spn{\ple{X,\eqrelr}}{[f]}{\ple{A,\rid_A}}{[g]}{\ple{Y,\eqrels}}$ in $\Span{\ct{C_{ex/wlex}}}_{\ct{C}_\Delta}(\ple{X,\eqrelr},\ple{Y,\eqrels})$, a span $\spn{X}{f}{A}{g}{Y}$ and use \cref{l:chiudo} to map $\gr{f}\rconv\rcomp\gr{g}$ in 
$\Span\CC(X,Y)$ to $\eqrelr\rcomp\gr{f}\rconv\rcomp\gr{g}\rcomp\eqrels$ in $\EQR{\Span{\CC}}$
$(\ple{X,\eqrelr},\ple{Y,\eqrels})$. This assignment does not depend on the choice of the representative: if  $\spn{X}{f'}{A}{g'}{Y}$ is another choice, then $\rid_A\order\gr{f}\rcomp\eqrelr\rcomp\gr{f'}\rconv$ and $\rid_A\order\gr{g'}\rcomp\eqrels\rcomp\gr{g}\rconv$ in $\Span{\CC}(A,A)$; equivalently, $\gr{f}\rconv\order\eqrelr\rcomp\gr{f'}\rconv$ and $\gr{g}\order\gr{g'}\rcomp\eqrels$, so
\[
\eqrelr\rcomp\gr{f}\rconv\rcomp\gr{g}\rcomp\eqrels
\order
\eqrelr\rcomp\eqrelr\rcomp\gr{f'}\rconv\rcomp\gr{g'}\rcomp\eqrels\rcomp\eqrels
=
\eqrelr\rcomp\gr{f'}\rconv\rcomp\gr{g'}\rcomp\eqrels
\]
analogously one shows that $\eqrelr\rcomp\gr{f'}\rconv\rcomp\gr{g'}\rcomp\eqrels\order\eqrelr\rcomp\gr{f}\rconv\rcomp\gr{g}\rcomp\eqrels$. This proves that $i'$ is well defined function.

 It is also monotone: take a bigger span $\spn{\ple{X,\eqrelr}}{[h]}{\ple{B,\rid_B}}{[l]}{\ple{Y,\eqrels}}$, i.e. suppose that there is $\fun{k}{A}{B}$ with $[hk]=[f]$ and $[lk]=[g]$. In terms of the relational operations of $\Span{\CC}$ it is $\gr{f}\rconv\order \eqrelr\rcomp\gr{h}\rcomp\gr{k}\rconv$
 and $\gr{g}\order \gr{k}\rcomp\gr{l}\rcomp\eqrels$. So
 \[
 \eqrelr\rcomp\gr{f}\rconv\rcomp\gr{g}\rcomp\eqrels
 \order
 \eqrelr\rcomp\eqrelr\rcomp\gr{h}\rcomp\gr{k}\rconv\rcomp\gr{k}\rcomp\gr{l}\rcomp\eqrels\rcomp\eqrels
 \order
 \eqrelr\rcomp\gr{h}\rcomp\gr{l}\rcomp\eqrels
 \]
where we use transitivity of $\eqrelr$ and $\eqrels$ and functionality of $\gr{k}$. 
\cref{l:isofrecce} tell us that if we show that $i'$ is the inverse of $i$ as a monotone function, then $i'$ commutes with relational operations, providing the inverse of $i$ in $\RDtn$.

Since elements of $\EQR{\Span\CC}(\ple{X,\eqrelr},\ple{Y,\eqrels})$ are descent data, the composition $i'i$ is the identity. Moreover, by reflexivity of $\eqrelr$ and $\eqrels$, it holds that  $\relr\order \eqrelr\rcomp\relr\rcomp\eqrels$, so $\id\order ii'$. It remains to show that $ii'\order\id$.
 Recall that $\eqrelr\rcomp\gr{f}\rconv\rcomp\gr{g}\rcomp\eqrels$ is given by the two external legs of the following diagram of \CC

\[
\vcenter{\xymatrix@C=3ex@R=3ex{
&&& W \ar[ld]_-{w_1} \ar[rd]^-{w_2} \ar@{}[dd]|{wpb} &&& \\
&& T\ar[ld]_-{t_1} \ar[rd]|{t_2} \ar@{}[dd]|{wpb} && V\ar[ld]|{v_1} \ar[rd]^-{v_2} \ar@{}[dd]|{wpb}&& \\  
&R\ar[ld]_-{\eqrelr_1}\ar[rd]_-{\eqrelr_2} && A \ar[ld]^-{f} \ar[rd]_-{g} && S \ar[ld]^-{\eqrels_1}\ar[rd]^-{\eqrels_2}& \\ 
X&& X && Y &&Y
}}
\]
To show that $i(\eqrelr\rcomp\gr{f}\rconv\rcomp\gr{g}\rcomp\eqrels)\order i(\gr{f}\rconv\rcomp\gr{g})$ we need to find $\fun{k}{W}{A}$ such that $[\eqrelr_1t_1w_1]=[fk]$ and $[gk]=[\eqrels_2v_2w_2]$ in $\CC_{ex/wlex}$. That is, we need to find an arrow $\fun{r}{W}{R}$   with $\eqrelr_1r=\eqrelr_1t_1w_1$ and $\eqrelr_2r=fk$
and an arrow  $\fun{s}{W}{S}$ with 
$\eqrels_1s=gk$ and $\eqrels_2s=\eqrels_2v_2w_2$. It is immediate to verify that 
\[
\xymatrix{W\ar[r]^-{v_1w_2}&A}\quad \xymatrix{W\ar[r]^-{t_1w_1}&R}
\quad \xymatrix{W\ar[r]^-{v_2w_2}&S}
\]
are the desired $k$,
 $r$ and $s$, respectively.
